% Source: https://artofproblemsolving.com/community/c3h2508740_what_is_totient_of_1
% Stable identity: aops-topic-2508740-q01
\textbf{(a)} The only integer $k$ with $1\le k\le1$ is $k=1$, and $\gcd(1,1)=1$. So $\varphi(1)=\boxed{1}$.

\textbf{(b)} There are no positive integers less than $1$, so that wording would give $0$ for $n=1$ (for $n\ge2$ the two definitions agree, since $\gcd(n,n)=n\ne1$). But $2$ and $1$ are relatively prime, so the multiplicative rule would give
\[\varphi(2)=\varphi(2\cdot1)=\varphi(2)\varphi(1)=\varphi(2)\cdot0=0,\]
while $\varphi(2)=1$ (only $k=1$ is coprime to $2$). So the rule requires $\varphi(1)=1$, which is why the definition uses $1\le k\le n$.
